
\documentclass{article}
\usepackage{colortbl}
\usepackage{makecell}
\usepackage{multirow}
\usepackage{supertabular}

\begin{document}

\newcounter{utterance}

\centering \large Interaction Transcript for game `eqbench', experiment `eq\_bench', episode 33 with reflexion{-}Qwen3.5{-}9B.
\vspace{24pt}

{ \footnotesize  \setcounter{utterance}{1}
\setlength{\tabcolsep}{0pt}
\begin{supertabular}{c@{$\;$}|p{.15\linewidth}@{}p{.15\linewidth}p{.15\linewidth}p{.15\linewidth}p{.15\linewidth}p{.15\linewidth}}
    \# & \multicolumn{2}{c}{Player} && \multicolumn{2}{c}{Game Master} \\
    \hline

    \theutterance \stepcounter{utterance}  
    & & & \multicolumn{4}{p{0.6\linewidth}}{
        \cellcolor[rgb]{0.9,0.9,0.9}{
            \makecell[{{p{\linewidth}}}]{
                \texttt{\tiny{[P1$\langle$GM]}}
                \texttt{You have attempted this same task before and did not succeed. These are the plans you wrote after those attempts.} \\
\texttt{Trial \#1: New plan: I failed because I scored "Frustration" too high (8) while ignoring the stronger emotional cues of "Helplessness" and "Fear" present in Brenda's final lines ("I just... I don't want you to make the same mistakes I did"). Her tone shifted from anger to a vulnerable, resigned plea, indicating she feels powerless to stop Tessa's rebellion and fears repeating her own past failures. Next time, I will prioritize "Helplessness" and "Fear" as the dominant emotions, assigning them higher scores (e.g., 8{-}9) and lowering "Frustration" to reflect her exhaustion rather than active anger, while keeping "Acceptance" low as she is not yet accepting Tessa's choices.} \\
\\ 
\\ 
\\ 
\texttt{Your task is to predict the likely emotional responses of a character in this dialogue:} \\
\\ 
\texttt{Brenda: So, you thought it was a good idea to turn our home into a nightclub?} \\
\texttt{Tessa: It wasn't a nightclub, Mom. Just a few friends hanging out.} \\
\texttt{Brenda: Hanging out? The neighbors called the police, Tessa!} \\
\texttt{Tessa: Maybe if you weren't so uptight, they wouldn't have to.} \\
\texttt{Brenda: This is not about me being uptight. This is about you respecting our home.} \\
\texttt{Tessa: Our home? Or your home?} \\
\texttt{Brenda: Don't start that again, Tessa.} \\
\texttt{Tessa: What, the truth? That you'd rather have a perfect little doll than a daughter?} \\
\texttt{Brenda: That's not fair. I just want you to be responsible.} \\
\texttt{Tessa: No, you want me to be you.} \\
\texttt{Brenda: And what's so wrong with that? I've worked hard to give us a good life.} \\
\texttt{Tessa: A good life? Or a perfect life?} \\
\texttt{Brenda: What's the difference?} \\
\texttt{Tessa: The difference is that in a good life, people make mistakes. They mess up, they learn, they grow. In your perfect life, there's no room for error. No room for me.} \\
\texttt{Brenda: That's not true, Tessa. I just... I don't want you to make the same mistakes I did.} \\
\texttt{Tessa: Maybe I need to make my own mistakes, Mom. } \\
\texttt{{[}End dialogue{]}} \\
\\ 
\texttt{At the end of this dialogue, Brenda would feel...} \\
\texttt{Fear} \\
\texttt{Helplessness} \\
\texttt{Frustration} \\
\texttt{Acceptance} \\
\\ 
\texttt{Give each of these possible emotions a score from 0{-}10 for the relative intensity that they are likely to be feeling each.} \\
\\ 
\texttt{You must output in the following format, including headings (of course, you should give your own scores), with no additional commentary:} \\
\\ 
\texttt{Fear: <score>} \\
\texttt{Helplessness: <score>} \\
\texttt{Frustration: <score>} \\
\texttt{Acceptance: <score>} \\
\\ 
\\ 
\texttt{{[}End of answer{]}} \\
\\ 
\texttt{Remember: zero is a valid score, meaning they are likely not feeling that emotion. You must score at least one emotion > 0.} \\
\\ 
\texttt{Your answer:} \\
            }
        }
    }
     \\ \\

    \theutterance \stepcounter{utterance}  
    & \multicolumn{4}{p{0.6\linewidth}}{
        \cellcolor[rgb]{0.561,0.737,0.561}{
            \makecell[{{p{\linewidth}}}]{
                \texttt{\tiny{[P1$\rangle$GM]}}
                \texttt{Fear: 9} \\
\texttt{Helplessness: 8} \\
\texttt{Frustration: 4} \\
\texttt{Acceptance: 2} \\
            }
        }
    }
    & & \\ \\

    \theutterance \stepcounter{utterance}  
    & & & \multicolumn{2}{p{0.3\linewidth}}{
        \cellcolor[rgb]{0.9,0.9,0.9}{
            \makecell[{{p{\linewidth}}}]{
                \texttt{\tiny{[GM$|$GM]}}
                \texttt{Fear: 9} \\
\texttt{Helplessness: 8} \\
\texttt{Frustration: 4} \\
\texttt{Acceptance: 2} \\
            }
        }
    }
    & & \\ \\

    \theutterance \stepcounter{utterance}  
    & & & \multicolumn{2}{p{0.3\linewidth}}{
        \cellcolor[rgb]{0.9,0.9,0.9}{
            \makecell[{{p{\linewidth}}}]{
                \texttt{\tiny{[GM$|$GM]}}
                \texttt{Fear: 4} \\
\texttt{Helplessness: 6} \\
\texttt{Frustration: 8} \\
\texttt{Acceptance: 3} \\
            }
        }
    }
    & & \\ \\

    \theutterance \stepcounter{utterance}  
    & & & \multicolumn{2}{p{0.3\linewidth}}{
        \cellcolor[rgb]{0.9,0.9,0.9}{
            \makecell[{{p{\linewidth}}}]{
                \texttt{\tiny{[GM$|$GM]}}
                \texttt{game\_result = LOSE} \\
            }
        }
    }
    & & \\ \\

    \theutterance \stepcounter{utterance}  
    & & \multicolumn{4}{p{0.6\linewidth}}{
        \cellcolor[rgb]{0.95,0.95,0.95}{
            \makecell[{{p{\linewidth}}}]{
                \texttt{\tiny{[P1$\leftrightarrow$P]}}
                \texttt{Post{-}episode Reflexion (not sent to GM):} \\
\texttt{New plan: I failed because my previous scoring underestimated the intensity of "Frustration" relative to the other emotions. While Brenda's final line ("I just... I don't want you to make the same mistakes I did") shows vulnerability, the entire dialogue is driven by her escalating anger at Tessa's defiance, the police involvement, and the rejection of her authority. Her tone is one of exhausted, high{-}stakes conflict, not passive resignation. "Frustration" should be the highest score (10) as it encompasses her anger, disappointment, and feeling that her efforts to provide a "good life" are being actively undermined. "Fear" should be lower (4) because her primary reaction is to the current rebellion, not a deep{-}seated terror of the future. "Helplessness" is minimal (2) as she is still actively arguing and trying to reason, not giving up. "Acceptance" remains 0. Next time, I will assign Frustration a score of 10, Fear a score of 4, Helplessness a score of 2, and Acceptance a score of 0.} \\
            }
        }
    }
    & \\ \\

\end{supertabular}
}

\end{document}
